% ============================================
% MEM-13
% Gerado de molambudos.md — build_miolo.py
% ============================================
\section*{\textbf{MEM-13} --- O Médico Novo}
\markboth{MEM-13 --- O Médico Novo}{}

\noindent \textit{Primeira sessão com Dr. Heitor Oliveira}

\noindent \textbf{1979 --- Primeira sessão com Dr. Heitor Oliveira}

\vspace{0.5em}\noindent\rule{\textwidth}{0.5pt}\vspace{0.5em}

\noindent Eu sabia que ele era novo no jeito de andar.

\noindent Os médicos velhos andam devagar. Eles já viram de tudo. Não têm pressa. Os novos andam rápido. Eles acham que vão descobrir alguma coisa que ninguém descobriu.

\noindent O Dr. Oliveira andava rápido.

\noindent Ele entrou na sala de sessões --- uma sala com duas cadeiras, uma mesa, uma janela gradeada --- e sentou na frente de mim. Colocou um bloco na mesa. Uma caneta. Ajeitou o jaleco.

\noindent \textit{\textquotedblleft{}Bom dia, Joaquim.\textquotedblright{}}

\noindent \textit{\textquotedblleft{}Bom dia.\textquotedblright{}}

\noindent Ele me olhou. Eu olhei ele.

\noindent \textit{\textquotedblleft{}Sou o Dr. Heitor Oliveira. Sou o novo psiquiatra da ala.\textquotedblright{}}

\noindent \textit{\textquotedblleft{}Eu sei.\textquotedblright{}}

\noindent \textit{\textquotedblleft{}Sabe como?\textquotedblright{}}

\noindent \textit{\textquotedblleft{}O cheiro.\textquotedblright{}}

\noindent Ele parou.

\noindent \textit{\textquotedblleft{}Cheiro?\textquotedblright{}}

\noindent \textit{\textquotedblleft{}Médico novo cheira a sabão. Médico velho cheira a naftalina.\textquotedblright{}}

\noindent Ele riu. Era uma risada nervosa.

\noindent \textit{\textquotedblleft{}Você é observador.\textquotedblright{}}

\noindent \textit{\textquotedblleft{}Sessenta e dois anos aqui dentro. A gente aprende.\textquotedblright{}}

\noindent Ele anotou alguma coisa.

\noindent \textit{\textquotedblleft{}Você sabe por que está aqui, Joaquim?\textquotedblright{}}

\noindent \textit{\textquotedblleft{}Porque matei três pessoas.\textquotedblright{}}

\noindent \textit{\textquotedblleft{}Por que matou?\textquotedblright{}}

\noindent \textit{\textquotedblleft{}Porque elas estavam no lugar errado na hora errada.\textquotedblright{}}

\noindent \textit{\textquotedblleft{}E por que elas estavam no lugar errado?\textquotedblright{}}

\noindent \textit{\textquotedblleft{}Porque eu decidi que era a hora.\textquotedblright{}}

\noindent Ele parou de anotar.

\noindent \textit{\textquotedblleft{}Você planejou os assassinatos?\textquotedblright{}}

\noindent \textit{\textquotedblleft{}Não precisei. A oportunidade apareceu. Eu aproveitei.\textquotedblright{}}

\noindent \textit{\textquotedblleft{}Você sente remorso?\textquotedblright{}}

\noindent Eu olhei para ele. O olho amarelo.

\noindent \textit{\textquotedblleft{}O senhor já sentiu fome, doutor?\textquotedblright{}}

\noindent \textit{\textquotedblleft{}Já.\textquotedblright{}}

\noindent \textit{\textquotedblleft{}Fome de verdade? Fome de comer terra, de comer a própria mão, de olhar pro outro e pensar 'ele tem carne'?\textquotedblright{}}

\noindent Ele não respondeu.

\noindent \textit{\textquotedblleft{}Quando a gente passa fome desse jeito, não sobra espaço pra remorso. O remorso é um luxo de quem tem o estômago cheio.\textquotedblright{}}

\noindent \textit{\textquotedblleft{}Você está dizendo que matou por fome?\textquotedblright{}}

\noindent \textit{\textquotedblleft{}Estou dizendo que a fome matou por mim. Eu fui só a ferramenta.\textquotedblright{}}

\noindent Ele anotou. A caneta dele escrevia rápido.

\noindent \textit{\textquotedblleft{}E essa fome ainda está aqui?\textquotedblright{}}

\noindent \textit{\textquotedblleft{}Nunca foi embora.\textquotedblright{}}

\noindent \textit{\textquotedblleft{}O que você sente agora?\textquotedblright{}}

\noindent Eu fechei os olhos.

\noindent \textit{\textquotedblleft{}O senhor quer a verdade?\textquotedblright{}}

\noindent \textit{\textquotedblleft{}Quero.\textquotedblright{}}

\noindent \textit{\textquotedblleft{}Eu sinto o cheiro do senhor.\textquotedblright{}}

\noindent Ele parou de escrever.

\noindent \textit{\textquotedblleft{}Que cheiro?\textquotedblright{}}

\noindent \textit{\textquotedblleft{}O senhor comeu carne no jantar. Carne de boi. O senhor passou molho na carne. Tinha cebola. O senhor bebeu vinho depois. Vinho tinto. O senhor escovou os dentes, mas o cheiro ainda está nos poros da sua pele.\textquotedblright{}}

\noindent Ele engoliu seco.

\noindent \textit{\textquotedblleft{}E o que mais?\textquotedblright{}}

\noindent \textit{\textquotedblleft{}O senhor tem medo. O senhor está com medo agora. O senhor está sentado numa cadeira tentando parecer profissional, mas o senhor quer levantar e ir embora. O senhor ouviu histórias sobre mim. O senhor acha que eu sou um monstro.\textquotedblright{}}

\noindent \textit{\textquotedblleft{}Você é?\textquotedblright{}}

\noindent Eu abri os olhos.

\noindent \textit{\textquotedblleft{}Todo mundo é, doutor. A diferença é que eu parei de esconder.\textquotedblright{}}

\vspace{0.5em}
\noindent\textit{\textbf{↪ Rotas:} \rota{DOC-04} \rotasep \rota{MEM-12} \rotasep \rota{MEM-14}}